% =============================================================================
%  Kookboek — Familie Spoor — Lulu wraparound cover (back + spine + front)
%  Build:  latexmk -xelatex -output-directory=. -jobname=KookboekFamilieSpoor-cover cover/cover.tex
%          (build.sh does this for you)
%
%  This is uploaded to Lulu separately from the interior file (main.tex).
%  It is a single page sized to trim + wrap on every outer edge, per
%  Lulu's hardcover casewrap cover requirements: https://help.lulu.com
%
%  This book is a HARDCOVER CASEWRAP. The case (this cover) is physically
%  LARGER than the interior's text-block trim (see kookboek.sty) — the
%  case wraps around the block with a small overhang on every outer edge,
%  which is normal for hardcover binding, so the two files are sized
%  independently on purpose; they don't need to (and shouldn't) match.
%
%  Unlike a paperback, Lulu does not publish a continuous spine-width
%  formula for hardcover casewrap covers — per Lulu's Developer Portal
%  ("How is spine width calculated?", help.lulu.com), casewrap spine width
%  comes from a fixed page-count bucket table (the same guide p14 table
%  scripts/lulu_lint.py checks against), with every page count in a bucket
%  getting the exact same spine width. Trim, wrap area and safety margin
%  below were read directly off a real template generated through Lulu's
%  Project Creation Tool, q6qdvpe-cover-template.pdf (originally for a
%  164-page interior). \spinew was confirmed against a freshly regenerated
%  template for a 223-page interior (0.813in / 20.65mm), which matches the
%  documented bucket value (223-250 pages -> 0.8125in / 20.6375mm) almost
%  exactly — confirming the bucket table, not per-page variation, is what
%  determines it. The interior is now 252 pages, which moved it into the
%  next bucket (251-278 pages -> 22mm per guide p14 / HARDCOVER_SPINE_TABLE);
%  \spinew is set to that table value directly since no real Project
%  Creation Tool template has been regenerated for a page count in this
%  bucket yet — regenerate one before ordering a real proof.
%
%  BEFORE ORDERING A REAL PROOF:
%   - Set \interiorpagecount below to the actual page count of the built
%     interior PDF (KookboekFamilieSpoor.pdf) if it changes again.
%   - If the new page count falls in a different guide p14 bucket, update
%     \spinew to that bucket's table value (see scripts/lulu_lint.py's
%     HARDCOVER_SPINE_TABLE, or the Lulu Developer Portal spine-width
%     article) — regenerating a template from Lulu's Project Creation Tool
%     is the ultimate check before a real order, but no longer strictly
%     needed just to know the number.
% =============================================================================
\documentclass[10pt]{article}

% ---- Trim, wrap area, safety margin, spine --------------------------------
% Hardcover casewrap case size + wrap area, per q6qdvpe-cover-template.pdf.
% Intentionally larger than kookboek.sty's interior trim (see note above).
\newlength\trimw  \setlength\trimw{195.33mm}
\newlength\trimh  \setlength\trimh{258.57mm}
\newlength\bleed  \setlength\bleed{15.87mm}   % Lulu wrap area: 0.625in (hardcover, not paperback bleed)
\newlength\safety \setlength\safety{15.87mm}  % Lulu safety margin from wrap edge: 0.625in

\newcommand{\interiorpagecount}{260} % <-- UPDATE to match the built interior PDF
\newlength\spinew
\setlength\spinew{22mm} % <-- Lulu guide p14 hardcover-spine bucket value for
% 251-278 pages (see scripts/lulu_lint.py's HARDCOVER_SPINE_TABLE). The
% interior moved from 250 to 252 pages, crossing into this new bucket from
% the previous 223-250 one. No real Project Creation Tool template has been
% regenerated for this bucket yet — do that before ordering a real proof.
% If the page count moves to a different bucket, update to that bucket's
% value (see scripts/lulu_lint.py's HARDCOVER_SPINE_TABLE).

% Back and front cover panels are each trim width plus one outer bleed edge.
\newlength\panelw \setlength\panelw{\dimexpr\bleed+\trimw\relax}
% x-offset (from the spread's bottom-left corner) of the spine's and front
% panel's left edge.
\newlength\spinex \setlength\spinex{\panelw}
\newlength\frontx \setlength\frontx{\dimexpr\panelw+\spinew\relax}

\newlength\spreadw \setlength\spreadw{\dimexpr 2\panelw+\spinew\relax}
\newlength\spreadh \setlength\spreadh{\dimexpr \trimh+2\bleed\relax}

% Lulu's own cover template (q6qdvpe-cover-template.pdf) shows the safety
% margin is NOT a uniform \safety inset from every panel edge. It's
% \bleed+\safety from each TRUE outer/wrap edge (top, bottom, and the
% panel's own outer side) — since the wrap area sits between that edge
% and the trim line, and safety is measured a further \safety in from
% there — but only \safety from each SPINE-facing edge, since the trim
% line already sits right at the spine with no wrap zone to cross.
% Confirmed against the template's guide rectangles directly (extracted
% with PyMuPDF): e.g. its front-panel safety box is x=[688.68,1152.36]pt,
% which is [\frontx+\safety, \frontx+\panelw-\bleed-\safety] to the point.
\newlength\marginsafety \setlength\marginsafety{\dimexpr\bleed+\safety\relax}
\newlength\boxw \setlength\boxw{\dimexpr\panelw-\safety-\marginsafety\relax}
\newlength\boxh \setlength\boxh{\dimexpr\spreadh-2\marginsafety\relax}
\newlength\boxbottom \setlength\boxbottom{\marginsafety}
\newlength\boxtop \setlength\boxtop{\dimexpr\spreadh-\marginsafety\relax}
\newlength\frontboxleft \setlength\frontboxleft{\dimexpr\frontx+\safety\relax}
\newlength\frontboxright \setlength\frontboxright{\dimexpr\frontx+\panelw-\marginsafety\relax}
\newlength\frontboxcenterx \setlength\frontboxcenterx{\dimexpr\frontboxleft+0.5\boxw\relax}
\newlength\backboxleft \setlength\backboxleft{\marginsafety}
\newlength\backboxright \setlength\backboxright{\dimexpr\panelw-\safety\relax}
\newlength\backboxcenterx \setlength\backboxcenterx{\dimexpr\backboxleft+0.5\boxw\relax}

\usepackage[paperwidth=\spreadw, paperheight=\spreadh, margin=0pt]{geometry}
\pagestyle{empty}

% ---- Fonts (kept in sync with kookboek.sty) --------------------------------
\usepackage{fontspec}
\setmainfont{EB Garamond}[
  Path = fonts/ , Extension = .ttf ,
  UprightFont = EBGaramond-Regular, ItalicFont = EBGaramond-Italic, BoldFont = EBGaramond-SemiBold,
  Ligatures = TeX ]
\newfontfamily\handfont{Nothing You Could Do}[
  Path = fonts/ , Extension = .ttf ,
  UprightFont = NothingYouCouldDo-Regular,
  Ligatures = TeX]
\newfontfamily\julianafont{EB Garamond}[
  Path = fonts/ , Extension = .ttf ,
  UprightFont = EBGaramond-Regular]

% ---- Colours (kept in sync with kookboek.sty) ------------------------------
\usepackage{xcolor}
\definecolor{paper}{HTML}{FFFFFF}
\definecolor{ink}{HTML}{000000}
\definecolor{terracotta}{HTML}{555555}
\definecolor{inkfaint}{HTML}{888888}

\usepackage{amssymb}
\usepackage{tikz}
\usepackage{graphicx}

% The vendored EB Garamond .ttf exposes no smcp/c2sc GSUB feature (see
% kookboek.sty's Numbers=OldStyle/Lining comment for the same root cause),
% so \scshape silently falls back to plain upright text instead of small
% caps. \labelcaps fakes the label look with tracked full capitals instead.
\newcommand{\labelcaps}[2]{\addfontfeatures{LetterSpace=#1}\MakeUppercase{#2}}

\begin{document}
\thispagestyle{empty}
\noindent
\begin{tikzpicture}[remember picture, overlay]
  % Whole-spread background.
  \fill[paper] (current page.south west) rectangle (current page.north east);

  % ------------------------------------------------------------ FRONT ------
  % Panel occupies x in [\frontx, \frontx+\panelw]. Border is sized to
  % \boxw x \boxh — Lulu's actual safety-margin box for this panel (see
  % the note by \marginsafety above) — instead of the panel's own
  % (asymmetric wrap/spine) proportions. Inset like the interior title
  % page (main.tex) — kept clear of every edge, including the spine.
  %
  % Instead of one large hero photo, the box is filled with a mosaic of
  % every recipe's hero/served-dish photo (the same images \heroimagefade
  % puts on each recipe's own page) — \frontgrid/\backgrid list every
  % \heroimagefade target in recipes/*.tex, split alphabetically 40/36 so
  % each panel's grid (5x8 front, 6x6 back — see \drawmosaic below) uses
  % every one of them with no leftovers and no repeats. Each tile is drawn
  % from cover/mosaic/<name>.png, not images/<name> directly:
  % scripts/generate_cover_mosaic_crops.py pre-crops a copy of each hero
  % image down to the smallest bounding box that contains the whole dish
  % (found the same border-connected-background flood fill
  % scripts/normalize_background.py uses, so the crop can only ever remove
  % background, never the dish itself) plus a small fixed margin — the
  % source images/ files are untouched, since they're also used at full
  % size by \heroimagefade in the interior book. Each list entry is
  % name/w/h: w and h are that *cropped* copy's own aspect ratio scaled to
  % fit inside its grid cell on the longer axis, like CSS's object-fit:
  % contain, with the shorter axis padded out to the cell in plain white
  % by \drawmosaic below — so the tile is cropped only through the
  % background that was trimmed off already, never through the dish, and
  % never distorted by a width=/height= stretch either.
  \def\frontgrid{aardappel-groente-vlees/32.718mm/20.985mm,arretjescake/32.718mm/21.733mm,asperge-lasagne/32.718mm/18.103mm,blinis-vier-toppings/32.718mm/17.391mm,bloemkoolpasta-hero/32.398mm/28.354mm,bloemkoolquiche/32.718mm/20.292mm,bobotie/32.718mm/15.106mm,broccoli-uit-de-oven-hero/32.718mm/20.445mm,broccolipasta-witvis-olijven-rode-ui/32.718mm/22.658mm,brood/32.718mm/17.707mm,broodje-hamburger/32.718mm/16.967mm,broodje-knakworst/32.718mm/19.319mm,burritos/32.718mm/18.654mm,chili-con-carne/32.718mm/24.710mm,ciabatta/32.718mm/21.845mm,couscous/32.718mm/25.108mm,doner-kebab/32.718mm/21.222mm,fajitas/30.888mm/28.354mm,falafel/32.718mm/18.390mm,flammkuchen/32.718mm/18.500mm,focaccia/32.718mm/19.042mm,gegrilde-courgette/32.718mm/25.905mm,geroerbakte-spruitjes-hero/32.718mm/24.894mm,gevulde-paprika-hero/32.718mm/21.975mm,gnocchi/32.718mm/18.338mm,gnocchi-sorrentina/32.718mm/18.151mm,gnocchi-zongedroogde-tomaten-hero/32.718mm/21.387mm,groene-salade/32.718mm/25.969mm,hamburgerbolletjes/32.718mm/20.488mm,horiatiki/30.789mm/28.354mm,involtini-di-maiale-hero/32.718mm/20.524mm,jambalaya/32.718mm/21.732mm,kalkoen-in-spek/32.718mm/19.924mm,kipfilet-parmezaan-knoflook/32.718mm/27.620mm,kippensoep-simpel/32.718mm/24.613mm,kofta/32.718mm/21.197mm,kokos-vis/32.718mm/24.852mm,lasagne-simpel/32.718mm/21.727mm,mac-n-cheese/32.718mm/21.499mm,meringues/32.718mm/22.255mm}
  \def\backgrid{nasi/27.265mm/24.475mm,paella-hero/27.265mm/18.499mm,parmigiana-melanzane/27.265mm/16.967mm,pasta-amatriciana/27.265mm/22.104mm,pasta-blauwe-kaas/27.265mm/19.911mm,pasta-bolognese/27.265mm/25.996mm,pasta-prei-hamblokjes-olijven/27.265mm/21.596mm,pasta-primavera/27.265mm/18.486mm,pasta-scarpariello/27.265mm/19.263mm,pastinaaksoep/27.265mm/22.201mm,pita/27.265mm/22.576mm,pretzels/27.265mm/13.705mm,quesadillas/27.265mm/16.518mm,quiche-chorizo-geitenkaas/27.265mm/14.568mm,quiche-groente-cirkel/27.265mm/19.801mm,quiche-groente-geitenkaas/27.265mm/16.144mm,quiche-ham-prei/27.265mm/15.152mm,quiche-lorraine/27.265mm/17.618mm,ravioli-rode-pesto-hero/27.265mm/19.343mm,roti/27.265mm/20.152mm,roze-koeken-hero/27.265mm/16.183mm,sajoer-boontjes-hero/27.265mm/20.278mm,salade-honing-mosterddressing-hero/27.265mm/20.761mm,saltimbocca/27.265mm/16.053mm,sauzijcenbroodjes/27.265mm/17.952mm,schiacciata-hero/27.265mm/14.240mm,semifreddo/27.265mm/19.181mm,shakshuka/27.265mm/17.836mm,spinazie-quiche/27.265mm/15.848mm,spruitenstamp-hero/27.265mm/20.372mm,tarte-tatin-witlof/27.265mm/16.791mm,tomaten-pestorisotto-hero/27.265mm/18.033mm,tompouce-geitenkaas-hero/27.265mm/20.966mm,tortellini-al-forno-hero/27.265mm/17.540mm,venkel-risotto-hero/27.265mm/29.249mm,worstjes-rode-rijst-salade/27.265mm/18.704mm}

  \newlength\frontcw \setlength\frontcw{\dimexpr\boxw/5\relax} % 5 columns
  \newlength\frontch \setlength\frontch{\dimexpr\boxh/8\relax} % 8 rows -> 40 tiles
  \newlength\backcw  \setlength\backcw{\dimexpr\boxw/6\relax}  % 6 columns
  \newlength\backch  \setlength\backch{\dimexpr\boxh/6\relax}  % 6 rows -> 36 tiles

  % \drawmosaic{name/w/h list}{columns}{cell width}{cell height}{panel box
  % left x}: tiles the list's images left-to-right, top-to-bottom into a
  % #2-wide grid filling the \boxw x \boxh box at x-offset #5. Each cell
  % is first painted paper-white, then the image is drawn at its own
  % pre-scaled, cell-fitting \w x \h (see above) centered in the cell.
  \newcommand{\drawmosaic}[5]{%
    \foreach \img/\w/\h [count=\gi from 0] in #1 {
      \pgfmathsetmacro{\gridrow}{floor(\gi/#2)}
      \pgfmathsetmacro{\gridcol}{\gi-#2*\gridrow}
      \pgfmathsetlengthmacro{\cellx}{#5+\gridcol*#3}
      \pgfmathsetlengthmacro{\celly}{\boxtop-\gridrow*#4-#4}
      \fill[paper] ([shift={(\cellx,\celly)}]current page.south west) rectangle
                   ([shift={(\dimexpr\cellx+#3\relax,\dimexpr\celly+#4\relax)}]current page.south west);
      \node[anchor=center, inner sep=0pt, outer sep=0pt] at
          ([shift={(\dimexpr\cellx+0.5#3\relax,\dimexpr\celly+0.5#4\relax)}]current page.south west) {%
        \includegraphics[width=\w, height=\h]{cover/mosaic/\img}};
    }
  }
  \drawmosaic{\frontgrid}{5}{\frontcw}{\frontch}{\frontboxleft}

  \draw[ink, line width=1.2pt]
    ([shift={(\frontboxleft,\boxbottom)}]current page.south west) rectangle
    ([shift={(\frontboxright,\boxtop)}]current page.south west);

  % Semi-opaque paper plates behind the title block and footer line so
  % they stay legible over the busy photo mosaic instead of sitting
  % directly on top of it. Snapped to whole grid-row boundaries (in units
  % of \frontch down from \boxtop) rather than an arbitrary \spreadh
  % fraction, so the plate edge always falls on a row seam — never through
  % the middle of a row, which would leave a sliced-off strip of a photo
  % peeking out above or below the text.
  \fill[paper, opacity=0.9, rounded corners=2pt]
    ([shift={(\dimexpr\frontboxleft+4mm\relax,\dimexpr\boxtop-4\frontch\relax)}]current page.south west) rectangle
    ([shift={(\dimexpr\frontboxright-4mm\relax,\dimexpr\boxtop-1\frontch\relax)}]current page.south west);
  \fill[paper, opacity=0.9, rounded corners=2pt]
    ([shift={(\dimexpr\frontboxleft+4mm\relax,\boxbottom)}]current page.south west) rectangle
    ([shift={(\dimexpr\frontboxright-4mm\relax,\dimexpr\boxtop-7\frontch\relax)}]current page.south west);

  % anchor=center (not south, as before the plates existed) so each text
  % block sits centred on its own plate regardless of how many lines it
  % has, instead of needing a hand-tuned y offset kept in sync by hand.
  \node[anchor=center, align=center, text width=\boxw] at
      ([shift={(\frontboxcenterx,\dimexpr\boxtop-2.5\frontch\relax)}]current page.south west) {%
    {\footnotesize\color{terracotta}\labelcaps{6.0}{Familierecepten \textperiodcentered{} bijeengebracht met zorg}}\par
    \vspace{2mm}
    {\fontsize{64}{64}\selectfont\addfontfeatures{LetterSpace=6.0}\color{ink}KOOKBOEK}\par
    \vspace{2mm}
    {\footnotesize\color{inkfaint}\labelcaps{6.0}{Recepten van de familie}}\par
    \vspace{1mm}
    {\julianafont\fontsize{40}{40}\selectfont\color{ink}Spoor}\par};
  \node[anchor=center, align=center] at
      ([shift={(\frontboxcenterx,\dimexpr\boxtop-7.5\frontch\relax)}]current page.south west) {%
    \footnotesize\color{ink}\labelcaps{6.0}{Bereid met liefde \textperiodcentered{} Anno MMXXVI}};

  % ------------------------------------------------------------- SPINE -----
  % Panel occupies x in [\spinex, \spinex+\spinew].
  % rotate=-90 (not 90) so the title reads right-side up when the finished
  % book lies flat on a table: with rotate=90 it came out upside down.
  \node[rotate=-90, anchor=center, align=center] at
      ([shift={(\spinex+0.5\spinew,0.5\spreadh)}]current page.south west) {%
    {\color{ink}\labelcaps{8.0}{Kookboek}}%
    \hspace{1.5em}\raisebox{0.1ex}{\textcolor{terracotta}{$\scriptstyle\blacklozenge$}}\hspace{1.5em}%
    {\julianafont\color{ink}Spoor}};

  % -------------------------------------------------------------- BACK -----
  % Panel occupies x in [0, \panelw]. Border is \boxw x \boxh again, same
  % as the front — Lulu's actual safety-margin box for this panel — but
  % centered at \backboxcenterx, not \frontboxcenterx: the box sits off
  % the panel's geometric center since the wrap-edge (left, here) margin
  % is bigger than the spine-edge (right) margin (see \marginsafety note).
  % Filled with the remaining 36 hero images (\backgrid, defined with
  % \frontgrid up in the FRONT block above) in the same tiled-mosaic style.
  \drawmosaic{\backgrid}{6}{\backcw}{\backch}{\backboxleft}

  \draw[ink, line width=1.2pt]
    ([shift={(\backboxleft,\boxbottom)}]current page.south west) rectangle
    ([shift={(\backboxright,\boxtop)}]current page.south west);

  % One paper plate behind the whole "Over dit boek" block (label, rule,
  % quote and signature) so it reads clearly over the mosaic. Snapped to
  % whole grid-row boundaries (in units of \backch down from \boxtop),
  % same reasoning as the front panel's plates: a plate edge landing
  % mid-row would leave a sliced-off strip of a photo showing through.
  \fill[paper, opacity=0.9, rounded corners=2pt]
    ([shift={(\dimexpr\backboxleft+4mm\relax,\dimexpr\boxtop-5\backch\relax)}]current page.south west) rectangle
    ([shift={(\dimexpr\backboxright-4mm\relax,\dimexpr\boxtop-1\backch\relax)}]current page.south west);

  \node[anchor=south] at
      ([shift={(\backboxcenterx,0.64\spreadh)}]current page.south west) {%
    \fontsize{8}{10}\selectfont\color{terracotta}\labelcaps{6.0}{Over dit boek}};
  \node[anchor=north] at
      ([shift={(\backboxcenterx,0.625\spreadh)}]current page.south west) {%
    \color{terracotta}\rule{2.4cm}{0.5pt}};

  \node[anchor=center, align=center, text width=\boxw] at
      ([shift={(\backboxcenterx,0.5\spreadh)}]current page.south west) {%
    \itshape\color{inkfaint}Recepten die ik jarenlang uit mijn hoofd kookte,\\
    eindelijk opgeschreven, hoeveelheden bij benadering.\\
    Proef vooral mee, en pas gerust aan waar nodig.};

  \node[anchor=north] at
      ([shift={(\backboxcenterx,0.5\spreadh-2.4cm)}]current page.south west) {%
    \raisebox{0.1ex}{\textcolor{terracotta}{$\scriptstyle\blacklozenge$}}};
  \node[anchor=north] at
      ([shift={(\backboxcenterx,0.5\spreadh-3.0cm)}]current page.south west) {%
    \handfont\fontsize{22}{22}\selectfont\color{ink}Ben Spoor};
\end{tikzpicture}
\end{document}
